\scanpage{58}

Soit $(SS)$ la catégorie des ensembles semi-simpliciaux
\[
(SS)=\widehat{\Delta}=\underline{\operatorname{Hom}}(\Delta^\circ,(\mathrm{Ens}))
\]
C'est un topos, et si
\[
X\in\operatorname{Ob}(SS)
\]
est un ens $\frac12$ simplicial, on regarde le topos induit
\[
(SS)_{/X}\simeq\widehat{\Delta/X}
\simeq\underline{\operatorname{Hom}}(\Delta^\circ\,\mathrm{Ens})
\]

Pour nous, l'étude cohomologique de $X$ est celle de ce topos induit. Les faisceaux \inserted{(disons)} ``sur $X$'' sont les \deleted{\unsure{gr}} faisceaux semi-simpliciaux \unsure{sur} $F\longrightarrow X$ sur $X$, ou encore les foncteurs $(\Delta/X)^\circ\longrightarrow(\mathrm{Ens})$; les faisceaux abéliens sur $X$ sont les objets groupe abélien $F\longrightarrow X$ sur $X$, ou encore les foncteurs $(\Delta/X)^\circ\mathbin{-}(\mathrm{Ab})$ etc. Idem pour faisceaux abéliens gradués différentiels etc. Si $A_*$ est un ens. ss. on peut regarder $A\times X\longrightarrow X$ comme faisceau sur $X$ (c'est aussi le faisceau induit par $A_*$ sur le topos induit $(SS/X)$), c'est aussi le foncteur composé
\diagram{\begin{tikzcd}
(\Delta^{\unsure{\circ}}/X)^\circ\arrow[r]&(\Delta)^\circ\arrow[r,"A"]&(\mathrm{Ens}).
\end{tikzcd}}

\margin{57}
